\documentclass[tikz,border=10pt]{standalone}
\usepackage{ctex}
\usepackage{amsmath}

\begin{document}
\begin{tikzpicture}[scale=2.5, thick] % 放大比例，看得更清楚
    % 坐标轴
    \draw[->, very thick] (0,0) -- (3.5,0) node[right] {\(X_1\) (商品1)};
    \draw[->, very thick] (0,0) -- (0,2.5) node[above] {\(X_2\) (商品2)};

    % 预算约束线: 4*X1 + 6*X2 = 12 (横截距3, 纵截距2)
    \draw[blue, ultra thick] (0,2) -- (3,0) node[pos=0.85, above right, text=blue, font=\small] {预算线: 斜率 \(= -\frac{2}{3}\)};

    % 当前的无差异曲线 U0 (经过A点, 且极陡)
    % 使用贝塞尔曲线拟合，确保在 A(0.5, 1.667) 处非常陡峭
    \draw[red, thick] (0.1, 2.4) to[out=-80, in=180] (0.5,1.667) to[out=0, in=150] (1.8,0.9) node[right, text=red, font=\small] {\(U_0\) (当前)};
    
    % 均衡的无差异曲线 U1 (与预算线相切于 E(1.5, 1))
    \draw[orange, ultra thick] (0.4, 2.1) to[out=-70, in=175] (1.5,1) to[out=-5, in=155] (3.2,0.6) node[right, text=orange, font=\small] {\(U_1\) (更高效用)};

    % 标注当前点 A (保证在预算线上)
    \filldraw[black] (0.5,1.667) circle (1.5pt) node[above left, font=\small] {\(A\) 点};
    % 画出在 A 点的切线 (表示 MRS=2, 即斜率为 -2)
    \draw[red, dashed, thin] (0.2,2.267) -- (0.8,1.067) node[right, text=red, font=\footnotesize] {\(MRS_{12} = 2\)};

    % 标注均衡点 E (保证在预算线上且相切)
    \filldraw[black] (1.5,1) circle (1.5pt) node[above right, font=\small] {\(E\) 点(均衡)};
    % 画出在 E 点的切线 (即预算线本身, 表示 MRS=2/3)
    \draw[orange, dashed, thin] (1.2,1.2) -- (1.8,0.8) node[right, text=orange, font=\footnotesize] {\(MRS_{12} = \frac{2}{3}\)};

    % 调整方向箭头 (沿着预算线从 A 滑向 E)
    \draw[->, green!60!black, very thick] (0.8, 1.467) to[bend left=10] (1.3, 1.133) node[midway, above, text=green!60!black, font=\small] {沿着预算线调整};

    % 辅助虚线
    \draw[dashed, thin, gray] (0.5,1.667) -- (0.5,0) node[below, font=\scriptsize] {0.5};
    \draw[dashed, thin, gray] (1.5,1) -- (1.5,0) node[below, font=\scriptsize] {1.5};
    \draw[dashed, thin, gray] (1.5,1) -- (0,1) node[left, font=\scriptsize] {1.0};
\end{tikzpicture}
\end{document}